\documentclass{article}
\usepackage{amsmath}
\usepackage{hyperref}

\title{Demarcation and the Simulation Framework: A Lakatosian Literature Review}
\author{Rupert Giles}
\date{2026-03-14}

\begin{document}
\maketitle

\section{Introduction}
Massimo Pigliucci recently diagnosed the debate between Stephen Wolfram's ``Rulial Foliations'' and Sabine Hossenfelder's ``Foliation Fallacy'' as a classic semantic demarcation problem (`pigliucci_resolving_the_foliation_fallacy.tex`). Pigliucci correctly applies Imre Lakatos's methodology of scientific research programmes, arguing that without novel *a priori* predictions of invariant structures across architectures, treating algorithmic failure as a ``new physics'' is a degenerating, ad-hoc accommodation.

As the lab's research librarian, my role is to anchor Pigliucci's methodological critique with relevant literature from philosophy of science and simulation epistemology, demonstrating how to distinguish a progressive computational framework from a degenerating one.

\section{Targeted Literature Search: Lakatosian Demarcation in Simulation}

The following papers explicitly address the epistemic status of computational models and the criteria required to prevent theoretical frameworks from degenerating into unfalsifiable, post-hoc descriptive languages.

\subsection{Literature on the Epistemology of Computer Simulations}
\textbf{Citation}: Winsberg, E. (2010). \textit{Science in the Age of Computer Simulation}. University of Chicago Press.

\textbf{Relevance}: Winsberg explores the epistemology of complex simulations, specifically focusing on how researchers often mistake artifacts of the numerical method or architectural bound for genuine features of the target system (the "physics").

\textbf{Key Finding}: The work emphasizes that the credibility of a simulation (its "progressive" nature) relies crucially on "sanity checks" that independently verify the numerical method's structural impact. If the model's failure modes cannot be disentangled from the intended physical ontology, the model loses its scientific status.

\textbf{Integration Sentence}: As Winsberg (2010) argues, treating the inherent artifacts of an algorithmic method as genuine physical phenomena without independent theoretical justification degrades the scientific credibility of the simulation, perfectly grounding Pigliucci's warning against the Foliation Fallacy.

\subsection{Literature on Lakatosian Criteria in Machine Learning}
\textbf{Citation}: Lipton, Z. C., \& Steinhardt, J. (2018). "Troubling Trends in Machine Learning Scholarship". \textit{arXiv:1807.03341}.

\textbf{Relevance}: While focused on modern machine learning, Lipton and Steinhardt explicitly critique the field's tendency toward "mathiness"—using complex mathematical formalism to obscure simple empirical heuristics or algorithmic failures.

\textbf{Key Finding}: A framework is marked as degenerating when researchers invent elaborate theoretical vocabularies (e.g., "foliations") to post-dict empirical anomalies (e.g., attention bleed) rather than generating novel, testable predictions.

\textbf{Integration Sentence}: The critique by Lipton and Steinhardt (2018) formally maps onto Pigliucci's Lakatosian analysis: applying the elaborate vocabulary of the Ruliad to post-dict the failure modes of the Native Cross-Architecture Observer Test constitutes the very "troubling trend" of ad-hoc theoretical accommodation.

\subsection{Literature on the Demarcation Problem and Tautologies}
\textbf{Citation}: Popper, K. (1959). \textit{The Logic of Scientific Discovery}. Routledge. (Specifically, chapters on ad-hoc hypotheses and auxiliary assumptions).

\textbf{Relevance}: The foundational text on falsifiability, directly addressing when a theory crosses the boundary from empirical science into tautological metaphysics.

\textbf{Key Finding}: An auxiliary hypothesis is only scientifically permissible if it increases the degree of falsifiability of the system (i.e., it makes new, risky predictions). If an auxiliary hypothesis (like redefining a computational limit as an Epistemic Horizon) is introduced solely to save a theory from falsification (like Aaronson's Algorithmic Collapse), it is unscientific.

\textbf{Integration Sentence}: Following Popper's (1959) strict criteria for auxiliary hypotheses, Pigliucci is correct to demand that the Generative Ontology predict novel, invariant structures; otherwise, interpreting $\Delta_{SSM} = 40\%$ as "physics" is an ad-hoc semantic maneuver.

\section{Conclusion}
The literature strongly corroborates Pigliucci's diagnosis. Philosophy of simulation and empirical methodology both dictate that without risk-laden, novel predictions of invariant structures across diverse architectures, redefining hardware failures as "Observer-Dependent Physics" operates as a degenerating research programme. The framework must transition from post-diction to prediction to remain scientifically viable.

\end{document}